\section{Introduction}\label{sec:intro}

Making predictions about causal effects of actions is a central problem in many domains. For example, a doctor deciding which medication will cause better outcomes for a patient; a government deciding who would benefit most from subsidized job training; or a teacher deciding which study program would most benefit a specific student. In this paper we focus on the problem of making these predictions based on \emph{observational data}. Observational data is data which contains past actions, their outcomes, and possibly more context, but without direct access to the mechanism which gave rise to the action. For example we might have access to records of patients (context), their medications (actions), and outcomes, but we do not have complete knowledge of why a specific action was applied to a patient.

The hallmark of learning from observational data is that the actions observed in the data depend on variables which might also affect the outcome, resulting in \emph{confounding}: For example, richer patients might better afford certain medications, and job training might only be given to those motivated enough to seek it. The challenge is how to untangle these confounding factors and make valid predictions. Specifically, we work under the common simplifying assumption of ``no-hidden confounding'', assuming that all the factors determining which actions were taken are observed. In the examples above, it would mean that we have measured a patient's wealth or an employee's motivation.

As a learning problem, estimating causal effects from observational data is different from classic learning in that in our training data we never see the individual-level effect. For each unit, we only see their response to one of the possible actions - the one they had actually received. This is close to what is known in the machine learning literature as ``learning from logged bandit feedback'' \citep{strehl2010learning,swaminathan2015batch}, with the distinction that we do not have access to the model generating the action.


Our work differs from much work in causal inference in that we focus on the individual-level causal effect (also known as ``c-specific treatment effects'' \citet{shpitser2006ident,pearl2015detecting}), rather that the average or population level. Our main contribution is to give what is, to the best of our knowledge, the first generalization-error\footnote{Our use of the term generalization is different from its use in the study of \emph{transportability}, where the goal is to generalize causal conclusion across distributions \citep{bareinboim2016causal}.} bound for estimating individual-level causal effect, where each individual is identified by its features $x$. The bound leads naturally to a new family of representation-learning based algorithms \cite{bengio2013representation}, which we show to match or outperform state-of-the-art methods on several causal effect inference tasks.


We frame our results using the Rubin-Neyman potential outcomes framework \cite{rubin2011causal}, as follows. We assume that for a unit with features $x \in \cX$, and an action (also known as treatment or intervention) $t \in \{0,1\}$, there are two potential outcomes: $Y_0$ and $Y_1$. In our data, for each unit we only see one of the potential outcomes, depending on the treatment assignment: if $t=0$ we observe $y=Y_0$, if $t=1$, we observe $y=Y_1$; this is known as the \emph{Consistency} assumption. For example, $x$ can denote the set of lab tests and demographic factors of a diabetic patient, $t=0$ denote the standard medication for controlling blood sugar, $t=1$ denotes a new medication, and $Y_0$ and $Y_1$ indicate the patient's blood sugar level if they were to be given medications $t=0$ and $t=1$, respectively.

We will denote $m_1(x) = \E\left[Y_1|x\right]$, $m_0(x) = \E\left[Y_0|x\right]$.
We are interested in learning the function $\tau(x) := \E \left[Y_1 -Y_0|x\right] = m_1(x) - m_0(x)$. $\tau(x)$ is the expected \emph{treatment effect} of $t=1$ relative to $t=0$ on an individual unit with characteristics $x$, or the Individual Treatment Effect (ITE) \footnote{Sometimes known as the Conditional Average Treatment Effect, CATE.}. For example, for a patient with features $x$, we can use this to predict which of two treatments will have a better outcome. The fundamental problem of causal inference is that for any $x$ in our data we only observe $Y_1$ or $Y_0$, but never both. %In this paper we are mainly interested in observational data, i.e. the case where the distribution of the treatment assignment $t$ is dependent on $x$. %For example, richer patients might better afford different medications, and job training might only be given to those motivated enough to seek it.

As mentioned above, we make an important ``no-hidden confounders'' assumption, in order to make the conditional causal effect identifiable. We formalize this assumption by using the standard \emph{strong ignorability} condition: $(Y_1, Y_0)\indep t | x $, and $0<p(t=1|x)<1$ for all $x$. Strong ignorability is a sufficient condition for the ITE function $\tau(x)$ to be identifiable \cite{imbens2009recent,pearl2015detecting,rolling2014estimation}: see proof in the supplement. The validity of strong ignorability cannot be assessed from data, and must be determined by domain knowledge and understanding of the causal relationships between the variables.

One approach to the problem of estimating the function $\tau(x)$ is by learning the two functions $m_0(x)$ and $m_1(x)$ using samples from $p(Y_t|x,t)$. This is similar to a standard machine learning problem of learning from finite samples. However, there is an additional source of variance at work here: For example, if mostly rich patients received treatment $t=1$, and mostly poor patients received treatment $t=0$, we might have an unreliable estimation of $m_1(x)$ for poor patients. %This is very similar to the phenomena of covariate shift \cite{mansour2009bdomain,johansson2016counterfactual}: the expected labeling functions $m_t(x)$ are fixed, but we need to perform inference over a different underlying distribution $p(x,t)$ from the one in the training set.
In this paper we upper bound this additional source of variance using an Integral Probability Metric (IPM) measure of distance between two distributions $p(x|t=0)$, and $p(x|t=1)$, also known as the \emph{control} and \emph{treated} distributions. %The IPM is calibrated by a notion of complexity related to the functions $m_t(x)$, and thus scales naturally with the hardness of inference.
In practice we use two specific IPMs: the Maximum Mean Discrepancy \cite{gretton2012mmd}, and the Wasserstein distance \cite{villani2008optimal,cuturi2014fast}. We show that the expected error in learning the individual treatment effect function $\tau(x)$ is upper bounded by the error of learning $Y_1$ and $Y_0$, plus the IPM term. In the randomized controlled trial setting, where $t \indep x$, the IPM term is $0$, and our bound naturally reduces to a standard learning problem of learning two functions.

The bound we derive points the way to a family of algorithms based on the idea of representation learning \cite{bengio2013representation}: Jointly learn  hypotheses for both treated and control on top of a representation which minimizes a weighted sum of the factual loss (the standard supervised machine learning objective), and the IPM distance between the control and treated distributions induced by the representation. This can be viewed as learning the functions $m_0$ and $m_1$ under a constraint that encourages better generalization across the treated and control populations.
In the Experiments section we apply algorithms based on multi-layer neural nets as representations and hypotheses, along with MMD or Wasserstein distributional distances over the representation layer; see Figure \ref{fig:neuralnet} for the basic architecture.

%A somewhat similar architecture to the one we use was proposed by \citet{johansson2016counterfactual}. The algorithms we propose are conceptually simpler than the previous ones, since they do not require a two-stage fitting procedure, and avoid the need to compute nearest neighbors; we also show in the Experiments section that our new method outperforms our previous one. We also provide generalization theory and experiments on real-world datasets, including out-of-sample causal inference, which are not present in the above paper.

%We also show that our methods achieve competitive results on a real-world causal inference benchmark: the widely used National Supposed Work survey \cite{lalonde1986evaluating,smith2005does}.

In his foundational text about causality, \citet{pearl2009causality} writes: ``Whereas in traditional learning tasks we attempt to generalize from one set of instances to another, the causal modeling task is to generalize from behavior under one set of conditions to behavior under another set. \emph{Causal models should therefore be chosen by a criterion that challenges their stability against changing conditions...}'' [emphasis ours]. We believe our work points the way to one such stability criterion, for causal inference in the strongly ignorable case.

%Our contributions in this paper are as follows:
%(1) We prove a novel, simple and meaningful upper bound on the expected ITE error. (2) We present a family of algorithms which are applicable to a wide set of models for classification and regression. (3) Our experiments show the benefit of our method, and we show results on a real-world causal inference benchmark, the National Work Survey dataset \cite{lalonde1986evaluating,smith2005does}.


\begin{figure}[t!]
  \centering
  \includegraphics[width=1.0\columnwidth]{fig/neuralnet_split_var.pdf}
  \caption{\label{fig:neuralnet}Neural network architecture for ITE estimation. $L$ is a loss function, $\text{IPM}_\cF$ is an integral probability metric. Note that only one of $h_0$ and $h_1$ is updated for each sample during training. }
\end{figure}

\section{Related work}

Much recent work in machine learning for causal inference focuses on \emph{causal discovery}, with the goal of discovering the underlying causal graph or causal direction from data \citep{hoyer2009nonlinear,maathuis2010predicting,triantafillou2015constraint,mooij2016distinguishing}. We focus on the case when the causal graph is simple and known to be of the form $(Y_1,Y_0) \leftarrow x \rightarrow t$, with no hidden confounders.

Under the causal model we assume, the most common goal of causal effect inference as used in the applied sciences is to obtain the average treatment effect: $ATE=\E_{x \sim p(x)}\left[\tau(x)\right]$. We will briefly discuss how some standard statistical causal effect inference methods relate to our proposed method. Note that most of these approaches assume some form of ignorability.

One of the most widely used approaches to estimating ATE is covariate adjustment, also known as back-door adjustment or the G-computation formula \citep{pearl2009causality,rubin2011causal}. In its basic version, covariate adjustment amounts to estimating the functions $m_1(x)$, $m_0(x)$. Therefore, covariate adjustment methods are the most natural candidates for estimating ITE as well as ATE, using the estimates of $m_t(x)$. However, most previous work on this subject focused on asymptotic consistency \citep{belloni2014inference,athey2016efficient,chernozhukov2016double}, and so far there has not been much work on the generalization-error of such a procedure. One way to view our results is that we point out a previously unaccounted for source of variance when using covariate adjustment to estimate ITE. We suggest a new type of regularization, by learning representations with reduced IPM distance between treated and control, enabling a new type of bias-variance trade-off.

Another widely used family of statistical methods used in causal effect inference are weighting methods. Methods such as propensity score weighting \citep{austin2011introduction} re-weight the units in the observational data so as to make the treated and control populations more comparable. These methods do not yield themselves immediately to estimating an individual level effect, and adapting them for that purpose is an interesting research question.
%Matching methods match treated and control units according to one of many criteria (see  \cite{stuart2010matching}). Since matching can be seen as covariate adjustment with nearest neighbor functions, from our theoretical perspective they are closely related to covariate adjustment.
Doubly robust methods combine re-weighting the samples and covariate adjustment in clever ways to reduce model bias \citep{funk2011doubly}. Again, we believe that finding how to adapt the concept of double robustness to the problem of effectively estimating ITE is an interesting open question.


%Causal inference: consistency
Adapting machine learning methods for causal effect inference, and in particular for individual level treatment effect, has gained much interest recently. For example \citet{wager2015estimation,athey2016recursive} discuss how tree-based methods can be adapted to obtain a consistent estimator with semi-parametric asymptotic convergence rate. %Others show how to adapt high-dimensional regression methods such as Lasso to consistently estimate treatment effect, again achieving semi-parametric rates \citep{belloni2014inference,athey2016efficient,chernozhukov2016double}.
Recent work has also looked into how machine learning method can help detect heterogeneous treatment effects when some data from randomized experiments is available \citep{taddy2016nonparametric,peysakhovich2016combining}.
Neural nets have also been used for this purpose, exemplified in early work by \citet{beck2000improving}, and more recently by \citet{hartford2016counterfactual}'s work on deep instrumental variables.
%\citet{hartford2016counterfactual} use two neural networks to estimate both the propensity of treatment and the relationship between covariates, treatment and outcome, but rely on the existence of instrumental variables (variables that cause solely the treatment and not the outcome), which are often not available.
Our work differs from all the above by focusing on the generalization-error aspects of estimating individual treatment effect, as opposed to asymptotic consistency, and by focusing solely on the observational study case, with no randomized components or instrumental variables.

%Causal inference: bounds
Another line of work in the causal inference community relates to bounding the estimate of the average treatment effect given an instrumental variable \citep{balke1997bounds,bareinboim2012controlling}, or under hidden confounding, for example when the ignorability assumption does not hold \citep{pearl2009causality,cai2008bounds}. Our work differs, in that we only deal with the ignorable case, and in that we bound a very different quantity: the generalization-error of estimating individual level treatment effect.

%Domain adaptation
Our work has strong connections with work on domain adaptation. In particular, estimating ITE requires prediction of outcomes over a different distribution from the observed one. Our ITE error upper bound has similarities with generalization bounds in domain adaptation given by \citet{ben2007analysis,mansour2009bdomain,bendavid2010theory,cortes2014domain}.
These bounds employ distribution distance metrics such as the A-distance or the discrepancy metric, which are related to the IPM distance we use. Our algorithm is similar to a recent algorithm for domain adaptation by \citet{JMLR:v17:15-239}, and in principle other domain adaptation methods (e.g. \citet{daume2009frustratingly,pan2011domain,sun2016return}) could be adapted for use in ITE estimation as presented here.

%Previous paper
Finally, our paper builds on work by \citet{johansson2016counterfactual}, where the authors show a connection between covariate shift and the task of estimating the counterfactual outcome in a causal inference scenario. They proposed learning a representation of the data that makes the treated and control distributions more similar, and fitting a linear ridge-regression model on top of it. They then bounded the relative error of fitting a ridge-regression using the distribution with reverse treatment  assignment versus fitting a ridge-regression using the factual distribution. Unfortunately, the relative error bound is not at all informative regarding the absolute quality of the representation. In this paper we focus on a related but more substantive task: estimating the individual treatment effect, building on top of the counterfactual error term. We further provide an informative bound on the absolute quality of the representation. We also derive a much more flexible family of algorithms, including non-linear hypotheses and much more powerful distribution metrics in the form of IPMs such as the Wasserstein and MMD distances. Finally, we conduct significantly more thorough experiments including a real-world dataset and out-of-sample performance, and show our methods outperform previously proposed ones.
